\section{The Word in a Part}

\subsection{A dated God}

Christianity did not ease the scandal of particularity; it consummated it. ``And the Word was made flesh, and dwelt among us.''\endnote{John 1:14, Douay--Rheims, checked 24 July 2026.} Paul gives the event a timestamp---``when the fulness of the time was come, God sent his Son, made of a woman, made under the law''---and Luke gives it a bureaucratic setting, a tax census, ``a decree from Caesar Augustus,'' an enrollment ``first made by Cyrinus, the governor of Syria.''\endnote{Galatians 4:4--5 and Luke 2:1--2, Douay--Rheims, checked 24 July 2026. The article cites the census as Luke's deliberate anchoring of the nativity in datable public administration; the well-known historical questions about the Quirinius enrollment's date are outside its scope and are not adjudicated here.} The creeds would later fix the other end of the life with a governor's name: suffered under Pontius Pilate. No founding documents of any movement have ever been less interested in escaping contingency. Where Lessing saw accidental truths of history as religion's embarrassment, the evangelists nail their claim to the calendar as if daring the reader to check.

The claim, note carefully, is not the modest one that a great teacher arose in Galilee. It is that the eternal Word---the one through whom, the tradition had just finished saying, every singular thing receives being---assumed a particular human nature and lived a bounded human life. The fourth-century writers who defended this saw exactly where the pressure lay, because their pagan critics pressed it. Athanasius, answering Greeks who scoffed at the Word appearing in a body, refuses to blink at the metaphysics. The Word who took ``a body, and that of no different sort from ours'' was not thereby diminished or contained: ``He was not, as might be imagined, circumscribed in the body, nor, while present in the body, was He absent elsewhere.''\endnote{Athanasius, \emph{De incarnatione Verbi} 8 and 17, in Archibald Robertson's public-domain translation (Nicene and Post-Nicene Fathers, second series, vol.~4), checked 24 July 2026 in a full-text copy of the NPNF edition. Section 17 is the classical statement of what later theology discussed as the \emph{extra}: the incarnate Word's continuing cosmic presence and activity.} The Incarnation is not God squeezing into a part and vacating the whole. It is the one who fills the whole making himself knowable in a part.

Athanasius then turns the objectors' own cosmology against them. You grant, he says, that the Word pervades the universe; you should then ``not think it absurd that a single human body also should receive movement and light from Him''---for a man who admits the mind works in the whole body but balks at its working in the toe ``would be thought foolish.'' And to the demand for a grander instrument---``Why then did He not appear by means of other and nobler parts of creation, and use some nobler instrument, as the sun, or moon, or stars, or fire, or air, instead of man merely?''---he answers with the sentence that collapses the whole aesthetic of the objection: ``the Lord came not to make a display, but to heal and teach those who were suffering.''\endnote{Athanasius, \emph{De incarnatione Verbi} 42--43, Robertson translation, checked 24 July 2026. The quoted objection is Athanasius's own statement of his adversaries' question; ``He takes to Himself as an instrument a part of the whole, His human body,'' ibid.~43, gives the positive principle: known in the part because unrecognized in the whole.} A sun-god would have been a spectacle. Spectacles are for impressing the healthy. The sick need someone to arrive, at their scale, through their door. The choice of the small instrument is not a concession to human weakness incidental to the mission; it \emph{is} the mission's shape, because what had gone wrong---among all creatures, Athanasius notes, only man had gone wrong---was small, local, and fleshly, and the repair had to go where the wound was.

\subsection{Why the flesh had to be someone's}

Irenaeus, a century earlier, had already seen that the rescue fails unless the particularity is real all the way down. Against teachers for whom Christ's flesh was celestial or apparent, he asks the unanswerable housekeeping question about the Virgin: ``why did He come down into her if He were to take nothing of her?'' If Christ did not receive ``the substance of flesh from a human being,'' the exchange at the heart of the Gospel never occurs; for the Word of God ``did, through His transcendent love, become what we are, that He might bring us to be even what He is Himself.''\endnote{Irenaeus, \emph{Adversus haereses} III.22.1--2 and V, preface, in the public-domain Ante-Nicene Fathers translation (Roberts and Rambaut), checked at the New Advent transcription, 24 July 2026.} A generic incarnation---God assuming ``humanity'' in the abstract---would touch no one, because no one is humanity in the abstract.

Aquinas later gave that instinct its technical form, and it is worth pausing on, because it answers the objector's ``why one man?'' with unexpected precision. Could the Word have assumed human nature ``abstracted from all individuals''? No---outside individuals, human nature has no concrete being to assume; apart from them it exists only in an intellect.\endnote{Aquinas, \emph{Summa theologiae} III, q.~4, a.~4, English Dominican translation, checked 24 July 2026.} Could he then have assumed it in \emph{all} individuals at once? Again no: ``It was unfitting for human nature to be assumed by the Word in all its supposita,'' for that would abolish the many persons the nature exists to multiply, and would empty the Head--members relation by which the one Christ is ``the First-born of many brethren.''\endnote{Aquinas, \emph{Summa theologiae} III, q.~4, a.~5, English Dominican translation, checked 24 July 2026.} The alternatives to the scandal, in other words, are not nobler incarnations; they are non-events. If God is to become man at all, God becomes \emph{a} man. The particularity that offends is simply what reality costs.

And why should he become man at all? Here Aquinas argues fittingness, not necessity, and the distinction is part of the article's method: after the fact, one can see how the event befits God; no one could have deduced it beforehand. It belongs to the very nature of the good to communicate itself---\emph{pertinet autem ad rationem boni ut se aliis communicet}---and so it befits the highest good to communicate itself in the highest way, which occurs, Aquinas says quoting Augustine, when God so joins created nature to himself ``that one Person is made up of these three---the Word, a soul and flesh.''\endnote{Aquinas, \emph{Summa theologiae} III, q.~1, a.~1, corpus; English wording from the Dominican translation, Latin checked in the Corpus Thomisticum text, 24 July 2026. Arguments from fittingness (\emph{convenientia}) show harmony, not necessity: they answer ``why is this worthy of God?'' rather than ``why must this have happened?''---and this article uses them only in that role.} The maximal self-gift of the universal Good is not a universal broadcast. It is a union so close that the gift has a face.

\subsection{The ditch from the other side}

Now return to Lessing. His ditch is real, given his cargo: contingent facts can never total up to necessary truths. But the Gospel never proposed to carry that cargo. It does not offer history as a proof of geometry-like theses; it claims that the decisive good is itself an event---that ``all other historical events happen once, and then they pass away, swallowed up in the past,'' as the Catechism will put the contrast, while this one, uniquely, does not.\endnote{\emph{Catechism of the Catholic Church} 1085, official English text, checked at the Vatican website, 24 July 2026. The paragraph's own claim---that the Paschal mystery, unlike every other historical event, ``cannot remain only in the past''---is taken up in the following section.} What human beings needed saving from was never their ignorance of necessary truths, which the philosophers were already supplying at market rates. It was a condition: mortality, guilt, the unraveling of concrete lives in concrete bodies. A condition is not refuted; it is healed. And healing, unlike proof, is necessarily particular---it happens to patients, one by one, or not at all. Lessing standing at his ditch resembles a man who, informed that the doctor has arrived, objects that no visit, being a mere event, can demonstrate the axioms of medicine. The objection is exact and beside the point. The doctor has not come to demonstrate. ``The Lord came not to make a display, but to heal.''

The last irony the objection must digest is that this most particular of events is the ground of the most universal claim ever made about human beings. The Council's sentence is famous: ``For by His incarnation the Son of God has united Himself in some fashion with every man.'' He worked with human hands, thought with a human mind, loved with a human heart\endnote{Second Vatican Council, \emph{Gaudium et spes} 22, official English translation at the Vatican website, checked 24 July 2026. The paragraph's universality claim is the council's own; its precise theological modality (``in some fashion,'' \emph{quodammodo}) is preserved and not inflated here.}---and precisely by being irrevocably one of us, he is related to all of us, as no abstraction ever related to anyone. Universality was not sacrificed at Bethlehem. It was funded there. To the Greeks, foolishness; but the argument of these pages has been that ``the foolishness of God is wiser than men.''\endnote{1 Corinthians 1:25, Douay--Rheims, checked 24 July 2026.}
